% Paper title: TrajScale: Decoupled Mixed-Resolution Sampling for Efficient High-Resolution Video Generation

\section{Introduction}

\begin{figure*}[!t]
\centering
\includegraphics[width=0.82\textwidth]{figures/fig_teaser.png}
\caption{Real-frame comparison and TrajScale inference. (a) Content-aligned
Native-HR (estimated), Trilinear, CRLU-only, and TrajScale-Q outputs for
prompt 05, seed 9705, frame 41, with identical \(1.5\times\) crops. The
reference follows a complete 368p trajectory; transition outputs are 720p.
(b) Native-HR performs all 50 evaluations at high resolution, whereas
TrajScale-Q and TrajScale-E execute the first 40 and 45 evaluations at low
resolution and retain 10 and 5 high-resolution suffix evaluations.}
\label{fig:teaser}
\end{figure*}

Recent latent video diffusion models such as Wan generate semantically
coherent videos with realistic motion and fine details \citep{wan2025},
but their inference cost grows aggressively with spatial resolution,
video length, and sampling steps. In DiT-based models, every denoising
iteration processes the full spatiotemporal token sequence. Executing the
entire trajectory at the target resolution forces early steps---which mainly
establish coarse structure and macro-motion---to bear the full burden of
high-resolution computation, severely limiting real-time deployment.
Since diffusion sampling inherently distributes generation across stages
(global layout in early steps, high-frequency textures in later steps),
a promising paradigm is mixed-resolution inference: modeling global
content in a low-resolution latent space first, then switching to the
target resolution for a short refinement suffix. This spatial-scaling
approach \citep{lightx2v2026,zheng2026mrflow} is fundamentally orthogonal
and complementary to temporal step-distillation methods.

Yet, switching resolution within a diffusion trajectory is substantially more
complex than resizing a latent tensor. The transition is governed by three
coupled error sources: the residual gap to the converged low-resolution endpoint,
the discrepancy between clean latent spaces at different resolutions, and the
residual error left for high-resolution refinement. Directly regressing from the
handoff state to a high-resolution endpoint forces a single Joint
Trajectory--Scale Mapper (JTSM) to absorb timestep-dependent trajectory error,
VAE-specific scale mismatch, and underdetermined high-frequency content. Our
transition-step scans show that this entanglement attenuates fine detail,
motivating a decomposition into localized subproblems.

We therefore propose \textbf{TrajScale}, a decoupled mixed-resolution sampling
framework with three components activated at the resolution transition: a
LoRA-based \textbf{Endpoint Alignment Adapter (EAA)} that calibrates the current
clean prediction toward the converged low-resolution endpoint; a
\textbf{Cross-Resolution Latent Upsampler (CRLU)} that maps clean latents across
spatial scales; and \textbf{Target-Resolution Re-noising (TRR)}, a
parameter-free operation that reconstructs a schedule-consistent
high-resolution state for suffix refinement. EAA and CRLU are trained separately
through model-endogenous self-distillation using supervision generated by the
frozen target generator.

Our primary contributions are summarized as follows:
\begin{itemize}
\item \textbf{We propose TrajScale}, a lightweight mixed-resolution sampling framework that executes a low-resolution prefix followed by a short high-resolution suffix, yielding controllable quality--efficiency operating points (TrajScale-Q/E).
\item \textbf{We decouple the resolution transition} into endpoint alignment, cross-resolution clean-latent upsampling, and target-resolution re-noising through EAA, CRLU, and TRR. The decomposition isolates trajectory, scale, and refinement errors that are entangled in joint mapping.
\item \textbf{We develop model-endogenous dual supervision} without external paired videos or auxiliary teachers, and establish cached-prefix consistency for handoff-local adaptation. Experiments on standard and distilled Wan samplers validate spatial acceleration and compatibility with few-step sampling.
\end{itemize}

\section{Related Work}

\subsection{Latent Video Generation and Cascaded Upscaling}
Latent Diffusion Models compress the generative process into a compact space for scalable high-resolution video synthesis \citep{rombach2022ldm}. Traditional frameworks or recent systems like LTX-2 and LUVE typically employ cascaded pipelines, using independent spatial latent upsamplers or post-generation refinement modules to elevate resolution \citep{blattmann2023videoldm, hacohen2026ltx2, zhao2026luve}. Similarly, video super-resolution methods focus on restoring already decoded RGB outputs \citep{zhou2024upscale, xu2024videogigagan}.

Unlike approaches that treat upscaling as an isolated post-processing stage,
TrajScale inserts CRLU directly into the frozen Wan sampling trajectory. TRR
then reconstructs the target-resolution noisy state at the designated schedule
level, allowing the remaining Wan steps to synthesize high-frequency content
with the base model's native generative prior.

\subsection{Mixed-Resolution Diffusion Inference}
The changing-resolution inference in LightX2V utilizes fixed interpolation to adjust latent sizes mid-sampling \citep{lightx2v2026}. Concurrently, RALU demonstrates that training-free latent upscaling suffers from high-frequency aliasing and a resolution-dependent noise--timestep mismatch, designing a mixed-resolution scheme for Diffusion Transformers \citep{jeong2025ralu}. For image flow matching, MrFlow combines low-resolution structure generation with low-noise high-resolution refinement \citep{zheng2026mrflow}. TrajScale extends this paradigm to video by learning cross-resolution
clean-latent upsampling inside the original sampling process while explicitly
calibrating the residual low-resolution endpoint error before the resolution
transition.

\subsection{Few-Step Distillation and On-Policy Trajectory Learning}
To accelerate inference, temporal distillation methods minimize sampling steps via trajectory regression or consistency mapping \citep{luo2023lcm, yin2024dmd2, zhai2024mcm}. Recent works like Self-Forcing and AnyFlow further emphasize enforcing strict consistency between student training states and actual inference trajectories to mitigate distribution shift \citep{huang2025selfforcing, gu2026anyflow}.

TrajScale focuses on spatial rather than temporal computation, making it orthogonal to step reduction. Since EAA remains inactive prior to the transition step, its parameter updates do not alter the low-resolution prefix trajectory. Leveraging this local activation property, Section~\ref{sec:cached} establishes the consistency conditions that allow cached prefix states to be directly accessed during inference without training--inference distribution drift.

\section{Method}
\label{sec:method}

\subsection{Decoupled Mixed-Resolution Transition}
\label{sec:decoupled_transition}

Let \(v_\phi(x_k,k,c)\) denote a frozen Wan flow predictor, where
\(x_k\) is the latent state before the \(k\)-th denoising evaluation,
\(c\) is the text condition, and \(T\) is the total number of
evaluations. Superscripts \(L\) and \(H\) denote the low- and
target-resolution latent spaces, respectively. For a transition after
evaluation \(s\), \(x_s^L\) is the low-resolution state presented to the
transition operator. We use \(z_T^L\) to denote the endpoint obtained by
completing the native low-resolution trajectory under a fixed prompt,
random seed, initial noise, scheduler, and guidance configuration. We
further denote by \((z_0^L,z_0^H)\) a pair of clean Wan-VAE latents
encoded from low- and high-resolution versions of the same video.

A direct resolution transition could be learned with a single Joint
Trajectory--Scale Mapper (JTSM),
\[
M_\omega:(x_s^L,s,c)\mapsto z_0^H.
\]
Such a mapper, however, must simultaneously remove the residual error
between the transition-time prediction and the converged low-resolution
endpoint, translate between two VAE-specific latent distributions, and
infer high-frequency content that is not uniquely determined by the
low-resolution input. These responsibilities depend on different
variables: the first varies strongly with the transition step and
scheduler, the second is primarily a cross-scale representation problem,
and the third is better resolved by the pretrained generator's
high-resolution prior. Learning them through one regression target
therefore entangles trajectory completion, resolution conversion, and
detail synthesis.

TrajScale instead factorizes the transition into three domain-specific
operations:
\[
x_s^L
\xrightarrow{\mathrm{EAA}:v_{\phi+\Delta\theta_s}}
\widetilde z_s^L
\xrightarrow{U_\psi^{\mathrm{CRLU}}}
\widehat z_s^H
\xrightarrow{\mathrm{TRR}}
x_{s+1}^H.
\]
The Endpoint Alignment Adapter (EAA) corrects only the residual discrepancy
to the converged low-resolution endpoint. The Cross-Resolution Latent
Upsampler (CRLU) then maps the calibrated clean latent to the target
spatial scale. Finally, Target-Resolution Re-noising (TRR) reconstructs a
scheduler-compatible high-resolution state from which the frozen model
continues its remaining denoising evaluations. EAA and CRLU are trained
independently, while TRR is parameter free.

This factorization yields a design-oriented decomposition of transition
error. Let \(E_{\mathrm{endpoint}}(s)\) represent the residual
low-resolution endpoint discrepancy, \(E_{\mathrm{scale}}(s)\) the
cross-resolution clean-latent reconstruction error, and
\(E_{\mathrm{refine}}(s)\) the error left to the high-resolution suffix.
The transition can be viewed conceptually as
\[
E_{\mathrm{transition}}(s)
\approx
E_{\mathrm{endpoint}}(s)
+E_{\mathrm{scale}}(s)
+E_{\mathrm{refine}}(s).
\]
A later transition generally reduces the first term and moves the CRLU
input closer to its clean-latent training distribution, but it also
shortens the high-resolution suffix available to synthesize fine detail.
The transition step therefore controls a quality--efficiency trade-off,
whereas the three TrajScale operations isolate the error sources that a
monolithic JTSM must absorb jointly.

\begin{figure*}[!t]
\centering
\includegraphics[width=0.84\textwidth]{figures/fig_talh_overview.png}
\caption{Overview of TrajScale. Top: a low-resolution prefix is converted
into a high-resolution suffix through endpoint alignment, clean-latent
upsampling, and target-resolution re-noising. Bottom: both trainable
modules receive model-endogenous supervision from the same frozen Wan
generator. EAA pairs a cached transition state with the endpoint of the
corresponding native low-resolution trajectory, while CRLU pairs clean
latents encoded from a generated high-resolution video and its
RGB-downsampled counterpart.}
\label{fig:overview}
\end{figure*}

\subsection{Model-Endogenous Trajectory--Scale Learning}
\label{sec:model_endogenous_learning}

TrajScale learns trajectory-side calibration and scale-side reconstruction
from two complementary supervision streams generated entirely by the
frozen target model. It requires neither external paired videos,
pretrained super-resolution weights, nor an auxiliary teacher. The two
streams share the same generator and VAE but correspond to different
conditional distributions:
\[
\underbrace{(x_s^L,z_T^L)}_{\text{trajectory-side supervision}}
\quad\text{and}\quad
\underbrace{(z_0^L,z_0^H)}_{\text{scale-side supervision}}.
\]
The first teaches EAA how a transition-time prediction differs from the
endpoint of its own low-resolution trajectory. The second teaches CRLU
how the target generator represents the same video content at two spatial
scales. Separating these objectives prevents trajectory-dependent errors
from contaminating the cross-resolution mapping and avoids joint
backpropagation through the large frozen denoiser and the upsampler.

\subsubsection{Endpoint Alignment Adapter}
\label{sec:eaa}

\paragraph{Handoff-local endpoint supervision.}
For each training condition, we run the frozen model with a fixed text
prompt, random seed, initial noise, scheduler, and Classifier-Free
Guidance (CFG) configuration. We cache the low-resolution state \(x_s^L\)
immediately before the designated transition and continue the same native
low-resolution trajectory to obtain its converged endpoint \(z_T^L\).
The pair \((x_s^L,z_T^L)\) preserves the exact semantic and stochastic
conditions of the transition while isolating the residual error that
would otherwise be passed directly to CRLU.

\paragraph{Localized low-rank adaptation.}
EAA is implemented with LoRA \citep{hu2022lora} and is activated only at
the designated transition step. The base Wan parameters remain frozen.
For a linear layer \(W\), the adapted weight is
\[
W'=W+\frac{\alpha}{r}BA,
\]
where \(r\) is the LoRA rank. Low-rank updates are injected into the
\(q\), \(k\), \(v\), and \(o\) attention projections and the
\(\mathrm{ffn}.0\) and \(\mathrm{ffn}.2\) linear layers. Under the
flow-matching parameterization, EAA produces the calibrated clean-latent
estimate
\[
\widetilde z_s^L
=x_s^L-\sigma_s
v_{\phi+\Delta\theta_s}(x_s^L,s,c),
\]
where \(\sigma_s\) is the scheduler coefficient at the transition step.
We optimize
\[
\mathcal L_{\mathrm{EAA}}
=
\lVert\widetilde z_s^L-z_T^L\rVert_1
+0.1\lVert\widetilde z_s^L-z_T^L\rVert_2^2.
\]
For earlier transitions, where frame-wise residual errors are more likely
to propagate through the scale conversion, we additionally apply a
temporal residual alignment term with weight \(0.05\). Although the same
target can be expressed as a velocity label
\(v^*=(x_s^L-z_T^L)/\sigma_s\), all reported models are optimized in the
clean-latent domain. EAA performs no spatial upsampling; its sole role is
to move the transition state toward the native low-resolution endpoint
before CRLU is applied.

\paragraph{Cached-prefix consistency.}
\label{sec:cached}
Because EAA is inactive before step \(s\), its low-rank parameters satisfy
\(\Delta\theta_k=0\) for all \(k<s\). Under deterministic sampling and
matched text embeddings, initial noise, scheduler, CFG scale, and other
conditioning inputs, the adapted and frozen processes therefore use
identical transition functions throughout the prefix. By induction,
\[
x_k^{\mathrm{adapted}}=x_k^{\mathrm{base}},
\qquad \forall k\leq s.
\]
Consequently, a prefix state cached from the frozen model is exactly the
state consumed by EAA. This property supports offline pair construction
and controlled comparisons among multiple transition operators without
introducing prefix rollout mismatch; the same cached state can also be
reused whenever an identical prompt--seed prefix is evaluated repeatedly.
The equality requires transition-local activation, deterministic prefix
sampling, and matched schedulers, guidance, and conditioning. Extending
the adaptation to earlier or multiple steps would change subsequent
states and would require on-policy student rollouts or synchronized
teacher trajectories.

\subsubsection{Cross-Resolution Latent Upsampler}
\label{sec:crlu}

\paragraph{Generator-endogenous scale supervision.}
We first use the frozen Wan model to generate high-resolution reference
videos from predefined prompts and random seeds. Each video is spatially
downsampled in RGB space, after which the high- and low-resolution
versions are encoded independently by the same frozen Wan VAE. This
produces aligned clean-latent pairs \((z_0^L,z_0^H)\) with identical frame
counts, semantics, and motion trajectories. Performing the resolution
change before VAE encoding is important: CRLU learns the actual
cross-resolution correspondence induced jointly by the target generator,
RGB degradation, and VAE, rather than imitating an arbitrary interpolation
rule inside latent space.

\paragraph{Spatiotemporal latent upsampling.}
CRLU maps a 16-channel Wan latent to a target-resolution latent while
preserving the temporal dimension:
\[
U_\psi^{\mathrm{CRLU}}:
\mathbb R^{B\times16\times F\times h\times w}
\rightarrow
\mathbb R^{B\times16\times F\times H\times W}.
\]
Inspired by the latent upsampler in LTX-2
\citep{hacohen2026ltx2}, the network follows an
encoder--resampler--decoder design. A 3D convolution projects the input
into a 256-dimensional hidden representation, followed by four
spatiotemporal residual blocks. A spatial resampler changes only the
height and width, after which four additional 3D residual blocks and a
final 3D convolution reconstruct the 16-channel output. The 3D blocks
allow adjacent latent frames to contribute to the reconstruction without
changing the total number of frames.

We instantiate the spatial resampler according to the required latent-grid
ratio. For the \(480\times832\rightarrow720\times1248\) transition, the
latent grid changes from \(60\times104\) to \(90\times156\). CRLU first
expands the channel dimension by a factor of nine, applies a \(3\times\)
spatial PixelShuffle, and then performs stride-2 downsampling with a fixed
binomial blur kernel:
\[
60\times104
\xrightarrow{\times3}
180\times312
\xrightarrow{/2}
90\times156.
\]
For \(368\times640\rightarrow720\times1248\), the latent grid changes
from \(46\times80\) to \(90\times156\). A \(2\times\) PixelShuffle first
produces a \(92\times160\) grid, which is center-cropped to the exact
target size. These resolution-specific operators implement the desired
rational scale changes while the surrounding spatiotemporal network
learns the latent correspondence.

\paragraph{Cross-resolution objectives.}
Given
\(\widehat z_0^H=U_\psi^{\mathrm{CRLU}}(z_0^L)\), CRLU is trained with
\[
\mathcal L_{\mathrm{CRLU}}
=
\mathcal L_{\mathrm{rec}}
+\lambda_{\mathrm{content}}\mathcal L_{\mathrm{content}}
+\lambda_{\mathrm{temp}}\mathcal L_{\mathrm{temp}},
\]
where
\[
\mathcal L_{\mathrm{rec}}
=
\mathbb E\!\left[
\sqrt{(\widehat z_0^H-z_0^H)^2+\epsilon^2}
\right],
\]
and
\begin{align}
\mathcal L_{\mathrm{content}}
&=\lVert D(\widehat z_0^H)-z_0^L\rVert_1,\\
\mathcal L_{\mathrm{temp}}
&=\lVert\Delta_F\widehat z_0^H-\Delta_Fz_0^H\rVert_1.
\end{align}
Here, \(D\) maps the prediction back to the low-resolution grid, and
\(\Delta_F\) denotes first-order temporal differences between adjacent
latent frames. We use
\(\lambda_{\mathrm{content}}=0.2\),
\(\lambda_{\mathrm{temp}}=0.1\), and \(\epsilon=10^{-3}\). The three
terms respectively constrain target-latent reconstruction, preservation
of low-frequency content, and temporal evolution.

A low-resolution latent does not uniquely specify every high-frequency
texture in its high-resolution counterpart. CRLU is therefore not tasked
with deterministically hallucinating all missing detail. Instead, it
produces a structurally faithful and temporally stable clean latent on the
target grid. The remaining underdetermined frequencies are delegated to
the frozen high-resolution suffix, preserving a clear division of labor
between learned scale conversion and generative refinement.

\subsection{Schedule-Consistent High-Resolution Continuation}
\label{sec:hr_continuation}

After endpoint calibration and cross-resolution upsampling, TrajScale has
a target-grid clean latent
\[
\widehat z_s^H
=U_\psi^{\mathrm{CRLU}}(\widetilde z_s^L),
\]
but this clean estimate is not itself a valid intermediate state for the
next denoising evaluation. TRR reconstructs the state at the intended
scheduler level using the next-step coefficient \(\sigma_{s+1}\):
\[
x_{s+1}^H
=(1-\sigma_{s+1})\widehat z_s^H
+\sigma_{s+1}\epsilon^H,
\qquad
\epsilon^H\sim\mathcal N(0,I).
\]
The noise is sampled directly on the target-resolution grid. Unlike an
Interpolated-Flow Transition that resizes a low-resolution noisy state or
velocity field, TRR first completes the transition in the clean-latent
domain and then reconstructs a high-resolution state consistent with the
existing scheduler. It therefore avoids carrying a resolution-dependent
noise pattern across scales and supplies the frozen denoiser with the
state type it was originally trained to process.

Starting from \(x_{s+1}^H\), the remaining evaluations are performed by
the unmodified Wan denoiser in the high-resolution latent space. This
suffix refines the transferred structure and uses the base model's native
prior to synthesize fine textures and other spatial frequencies that are
not recoverable from \(z_0^L\) alone. TrajScale consequently assigns
coarse semantic layout and motion formation to the inexpensive
low-resolution prefix, cross-scale structure transfer to CRLU, and
high-frequency completion to a short high-resolution suffix.

The same transition rule applies to both conventional many-step samplers
and temporally distilled few-step samplers. Given any valid transition
index \(s\), the sampler executes its prefix at low resolution, invokes
EAA, CRLU, and TRR once, and completes the remaining scheduler states at
the target resolution. The framework does not alter the base scheduler or
require retraining the frozen denoiser. Concrete transition indices,
scheduler settings, and quality- or efficiency-oriented operating points
are specified in Section~\ref{sec:experiments}, since they are
instantiations of the general mixed-resolution procedure rather than part
of its definition.

\section{Experiments}
\label{sec:experiments}

\subsection{Experimental Setup}

We evaluate TrajScale at both system and module levels, covering end-to-end
quality--latency trade-offs, the individual and interactive effects of
CRLU and EAA, and compatibility with a 4-step distilled sampler.

For the 50-step pipeline, we generate approximately 1,000 reference
videos using Wan2.1 T2V-1.3B. Each video contains 81 frames at
\(720\times1248\) resolution and 16 fps. Bicubic RGB downsampling
produces corresponding \(480\times832\) and \(368\times640\) videos,
which are re-encoded by the frozen Wan VAE to construct aligned
cross-resolution latent pairs. EAA training pairs consist of the
low-resolution state immediately preceding a designated transition and the
endpoint of the corresponding complete low-resolution trajectory.

The 4-step experiments use approximately 5,000 videos generated by
Wan2.1 T2V-14B StepDistill-CfgDistill. CRLU is trained on
\(368\times640\rightarrow720\times1248\) clean latent pairs, while EAA
pairs the state preceding step 3 with the complete step-4
low-resolution endpoint. Training, validation, and test partitions are
separated by both text prompts and random seeds.

CRLU is optimized with AdamW for up to 50k steps using a learning rate of
\(1\times10^{-4}\), \(bf16\) precision, and an EMA decay of 0.9999.
EAA is trained for up to 10k steps with AdamW, a learning rate of
\(5\times10^{-5}\), and LoRA rank/scale \(32/32\). Training uses four
NVIDIA H100 GPUs and requires approximately 8 hours for CRLU and 33 hours
for EAA. All efficiency measurements are averaged over five independent
cold-start runs. Since all supervision is generated by the frozen target
model, TrajScale requires neither external paired videos nor an additional
upscaling teacher.

\subsection{Baselines and Evaluation Protocol}
\label{sec:factorialsetup}

System-level comparisons include \textbf{Native-HR Sampling}, which
executes all denoising steps at \(720\text{p}\);
\textbf{Trilinear Handoff}, which performs mixed-resolution inference
using fixed trilinear latent interpolation; and
\textbf{Endpoint Re-noising}, which completes the low-resolution
trajectory, applies CRLU, and performs one low-intensity
high-resolution refinement step. We additionally report
\textbf{Native-LR Sampling} as a low-resolution reference.

Module interactions are evaluated through a \(2\times2\) factorial
design that crosses handoff states
\(\{\text{Unaligned},\text{EAA-aligned}\}\) with resolution operators
\(\{\text{Trilinear},\text{CRLU}\}\). The four resulting configurations
are Trilinear Handoff, EAA+Trilinear, CRLU-only, and the complete TrajScale
framework. Oracle-Aligned CRLU and Joint Trajectory--Scale Mapper (JTSM)
are used as diagnostic references; extended results are provided in the
Supplementary Material.

For endpoint alignment, we report endpoint latent \(L_1\), decoded
PSNR and LPIPS, temporal latent differences, and sample-wise win rates.
CRLU is evaluated using latent \(L_1\), PSNR, SSIM, LPIPS, and temporal
\(L_1\). End-to-end quality is measured with VBench
\citep{huang2024vbench} and human evaluation. We define
\textbf{VBench-5} as the unweighted mean of subject consistency,
background consistency, motion smoothness, aesthetic quality, and
imaging quality.

All paired configurations share identical prompts, random seeds,
initial noise, schedulers, and guidance parameters. Prompt-level score
differences are analyzed using 10,000 paired bootstrap iterations for
\(95\%\) confidence intervals and two-sided exact sign tests. Human
evaluation uses randomized, double-blind comparisons over 10 matched
prompts, with three independent judges per prompt and prompt-level
majority voting. Inter-rater agreement is reported in the Supplementary
Material.

Efficiency measurements include end-to-end cold-start latency, peak GPU
memory, low- and high-resolution denoising evaluations, and speedup over
Native-HR Sampling. All EAA, CRLU, TRR, VAE, and dynamic model-loading
overheads are included.

\begin{figure*}[!t]
\centering
\includegraphics[width=0.76\textwidth]{figures/fig_quality_efficiency.png}
\caption{Quality--efficiency trade-offs relative to Native-HR Sampling.
(a) TrajScale-Q and TrajScale-E provide quality- and efficiency-oriented operating
points, while Endpoint Re-noising achieves higher speed at a larger
quality cost. (b) Absolute per-dimension VBench changes relative to
Native-HR.}
\label{fig:quality}
\end{figure*}

\begin{table*}[!t]
\centering
\footnotesize
\setlength{\tabcolsep}{5pt}
\begin{tabular}{lrrrrrr}
\toprule
Method & LR/HR evals. & Time (s) $\downarrow$ & Latency red. &
Speedup $\uparrow$ & VBench-5 $\uparrow$ & Retention $\uparrow$ \\
\midrule
Native-HR Sampling & 0/50 & $253.10\pm1.53$ & --- &
$1.00\times$ & \textbf{0.82836} & 100.00\% \\
TrajScale-Q (Step 40) & 40/10 & $138.36\pm0.50$ & 45.33\% &
$1.83\times$ & 0.80983 & 97.76\% \\
TrajScale-E (Step 45) & 45/5 & $\mathbf{114.26\pm1.52}$ &
\textbf{54.85\%} & $\mathbf{2.22\times}$ & 0.80792 & 97.53\% \\
\bottomrule
\end{tabular}
\caption{System-level comparison against Native-HR Sampling. Timings
average five cold-start runs, and quality is evaluated on ten matched
prompts with all processing overheads included.}
\label{tab:system}
\end{table*}

\section{Results and Analysis}

\subsection{System-Level Quality--Efficiency Trade-off}

Table~\ref{tab:system} compares the standard 50-step Native-HR pipeline
with the two principal TrajScale operating points. TrajScale-Q moves the first 40
evaluations to the low-resolution grid and retains 10 high-resolution
suffix evaluations, reducing latency from 253.10 to 138.36 seconds
(\(1.83\times\)). TrajScale-E retains only five high-resolution evaluations
and further reduces latency to 114.26 seconds (\(2.22\times\)).
Their VBench-5 scores are 0.80983 and 0.80792, compared with 0.82836 for
Native-HR Sampling.

The Endpoint Re-noising baseline reaches 86.45 seconds
(\(2.93\times\)) but decreases VBench-5 to 0.80093, with inferior
scores on \(9/10\) prompts relative to Native-HR
(\(p=0.02148\)). TrajScale therefore retains a longer high-resolution suffix
when additional structural and texture refinement is beneficial, while
later transitions prioritize inference speed. Peak memory remains
approximately 26 GiB across Native-HR, TrajScale-Q, TrajScale-E, and Endpoint
Re-noising, indicating that the primary system benefit of TrajScale is reduced
runtime rather than memory consumption.

\subsection{Component Analysis}

\paragraph{Cross-resolution latent upsampling.}
Table~\ref{tab:operator} evaluates CRLU against trilinear interpolation
over 50 paired clean latent samples. For
\(480\text{p}\rightarrow720\text{p}\), CRLU reduces latent \(L_1\) from
0.1929 to 0.0991, corresponding to a \(48.6\%\) reduction. It also
improves decoded PSNR by 5.71 dB, reduces LPIPS from 0.2358 to 0.0629,
and lowers temporal \(L_1\) by \(30.8\%\). CRLU wins on all 50 samples
for PSNR, LPIPS, and temporal consistency, and on 49/50 samples for
SSIM.

The more aggressive \(368\text{p}\rightarrow720\text{p}\) transition
shows the same trend: latent \(L_1\), LPIPS, and temporal \(L_1\) are
reduced by \(33.4\%\), \(51.1\%\), and \(15.5\%\), respectively.
These consistent latent, perceptual, and temporal improvements indicate
that CRLU learns a substantially more faithful cross-resolution mapping
than fixed interpolation.

\begin{table*}[!t]
\centering
\footnotesize
\setlength{\tabcolsep}{5pt}
\begin{tabular}{llrrrrr}
\toprule
Input $\rightarrow$ Output & Operator & Latent $L_1\downarrow$ &
PSNR $\uparrow$ & SSIM $\uparrow$ & LPIPS $\downarrow$ &
Temp. $L_1\downarrow$ \\
\midrule
480p $\rightarrow$ 720p & Trilinear & 0.1929 & 24.29 & 0.7159 &
0.2358 & 0.02121 \\
480p $\rightarrow$ 720p & CRLU & \textbf{0.0991} & \textbf{30.00} &
\textbf{0.8696} & \textbf{0.0629} & \textbf{0.01468} \\
368p $\rightarrow$ 720p & Trilinear & 0.2444 & 23.42 & 0.6740 &
0.3352 & 0.02388 \\
368p $\rightarrow$ 720p & CRLU & \textbf{0.1628} & \textbf{24.10} &
\textbf{0.7592} & \textbf{0.1640} & \textbf{0.02017} \\
\bottomrule
\end{tabular}
\caption{CRLU versus trilinear interpolation on paired clean latents
(\(n=50\)).}
\label{tab:operator}
\end{table*}

\paragraph{Endpoint alignment.}
At an alignment intensity of 0.75, EAA reduces the residual
low-resolution endpoint \(L_1\) error by \(25.82\%\) at step 40 and
\(21.03\%\) at step 45. Both settings improve all 10 evaluation
samples, with paired confidence intervals excluding zero. EAA also
increases decoded PSNR by 2.55 and 2.23 dB and reduces latent temporal
difference errors by \(23.44\%\) and \(21.24\%\), respectively. On the
4-step distilled model, EAA provides a smaller but consistent
\(5.03\%\) endpoint-error reduction. All three comparisons have a
10/0 paired win rate and two-sided exact sign-test \(p=0.00195\); full
confidence intervals are reported in Supplementary Table S2.

Alignment-intensity validation reveals a non-monotonic dependence on
correction strength. An intermediate intensity of 0.75 provides the
lowest endpoint error and the strongest downstream quality, whereas
weaker or full-strength corrections are less effective. This behavior
supports calibrated rather than unconstrained endpoint correction.

\paragraph{Factorial effects.}
The \(2\times2\) analysis in Table~\ref{tab:factorial} shows that CRLU
provides the principal broad quality gain, improving VBench-5 by
approximately \(0.031\)--\(0.041\) across the three evaluated
pipelines. EAA serves a more targeted role. With CRLU fixed, it improves
the imaging-quality dimension by \(+0.00599\), \(+0.00254\), and
\(+0.00479\) at step 40, step 45, and on the 4-step sampler,
respectively. These detail-oriented improvements are distributed across
only one component of the five-dimensional aggregate, and therefore
appear as smaller changes in the averaged VBench-5 score.

\begin{table*}[!t]
\centering
\footnotesize
\setlength{\tabcolsep}{4.5pt}
\begin{tabular}{lrrrrrrr}
\toprule
Sampler / transition & Trilinear & EAA+Tri. & CRLU-only & TrajScale &
CRLU gain & EAA gain & Total gain \\
\midrule
Wan50 / Step 40 (TrajScale-Q) & 0.77812 & 0.77901 & 0.80896 &
\textbf{0.80983} & +0.03085 & +0.00087 & +0.03171 \\
Wan50 / Step 45 (TrajScale-E) & 0.76776 & 0.76734 &
\textbf{0.80842} & 0.80792 & +0.04066 & -0.00050 & +0.04016 \\
Distill4 / Step 3-of-4 & 0.81796 & 0.81909 & 0.85670 &
\textbf{0.85680} & +0.03874 & +0.00010 & +0.03884 \\
\bottomrule
\end{tabular}
\caption{End-to-end factorial results on VBench-5 (\(n=10\)).
``EAA gain'' denotes the marginal effect of enabling EAA with CRLU
fixed.}
\label{tab:factorial}
\end{table*}

The results support a clear division of labor. CRLU restores the
cross-resolution latent representation and supplies the dominant
improvements in motion smoothness, aesthetics, and overall imaging
quality. EAA further calibrates the handoff state toward the converged
low-resolution endpoint, enriching localized textures and fine structures that may
be underrepresented by aggregate automatic metrics.

\subsection{Perceptual Effects and Distilled-Sampler Compatibility}

Human evaluation confirms that the detail recovery introduced by EAA is
perceptually visible. With the same CRLU operator, EAA receives
fine-detail win/loss/tie counts of \(8/2/0\) at step 40 and \(9/0/1\)
at step 45. Overall-quality judgments are also positive, reaching
\(6/2/2\) and \(6/0/4\), respectively. These preferences support the
dimension-level VBench observation that EAA selectively improves
imaging quality and localized structure.

The additional fine-scale content introduces a spatial--temporal
trade-off: temporal stability favors the unaligned CRLU state in both
settings. This suggests that EAA shifts the operating point toward
sharper local rendering, while CRLU and the remaining high-resolution
suffix provide the broader spatial and temporal reconstruction.

\begin{table}[!t]
\centering
\footnotesize
\setlength{\tabcolsep}{2pt}
\begin{tabular}{lrrr}
\toprule
Comparison & Details & Overall & Temporal \\
\midrule
EAA vs.\ Unaligned, CRLU (S40) & 8/2/0 & 6/2/2 & 1/9/0 \\
EAA vs.\ Unaligned, CRLU (S45) &
\textbf{9/0/1} & \textbf{6/0/4} & 0/8/2 \\
CRLU-only vs.\ Trilinear (S45) &
\textbf{10/0/0} & \textbf{10/0/0} & \textbf{8/0/2} \\
TrajScale-E vs.\ Trilinear (S45) &
\textbf{10/0/0} & \textbf{10/0/0} & \textbf{9/0/1} \\
\bottomrule
\end{tabular}
\caption{Prompt-level majority votes (W/L/T) over ten prompts with three
raters per prompt. Bold entries have exact sign-test \(p<0.05\).}
\label{tab:human}
\end{table}

CRLU provides the broader perceptual improvement: at step 45, both
CRLU-only and TrajScale-E obtain \(10/10\) prompt-level wins over trilinear
handoff for fine details and overall quality, while also improving
temporal stability. EAA adds a narrower refinement of local detail on
top of this cross-resolution reconstruction.

TrajScale is also compatible with the 4-step distilled Wan sampler. On the
\(368\text{p}\) trajectory, Trilinear Handoff achieves a VBench-5 score
of 0.81796, whereas TrajScale-D4 reaches 0.85680. The associated cold-start
latency increases only from 167.09 to 170.29 seconds. As in the 50-step
setting, CRLU supplies the primary cross-resolution quality improvement,
while EAA performs localized endpoint calibration and improves the
imaging-quality component.

Qualitative comparisons in Figure~\ref{fig:teaser}(a) further illustrate
CRLU's recovery of high-frequency textures and EAA's localized
adjustments. Extended spatial crops and temporal sequences are provided
in Figure S1. The content-matched
\texttt{Native-HR (estimated)} reference is produced from a complete
\(368\text{p}\) trajectory and is used only for frame-level structural
comparison; true \(720\text{p}\) Native-HR Sampling is used for the
quantitative evaluation in Table~\ref{tab:system}.

\section{Discussion and Limitations}

TrajScale assigns endpoint calibration, cross-resolution clean-latent
upsampling, and schedule-consistent state reconstruction to EAA, CRLU, and
TRR. This decomposition keeps CRLU in a stable clean-latent domain while the
frozen high-resolution suffix recovers underdetermined detail. Endpoint
metrics, imaging-quality gains, and human preferences indicate that EAA
improves localized structure and perceptual detail, although its effect depends
on the CRLU input distribution and the remaining suffix length.

EAA and CRLU use model-endogenous self-distillation: the frozen Wan model
supplies high-resolution references for CRLU and complete low-resolution
trajectories for EAA, avoiding external paired data, standalone
upscalers, and additional teachers.

Several limitations remain. CRLU is tied to the target VAE, generator
distribution, and scaling configuration, so cross-model transfer is
unverified. A distribution gap persists between its training inputs and
native low-resolution trajectory states. Fixed handoff steps also
require step-specific EAA parameters, motivating timestep-conditioned
adapters for adaptive handoffs. Finally, TrajScale addresses in-trajectory
spatial scaling rather than real-world degradations such as compression,
blur, or sensor noise. It inherits the base generator's content risks
and safety requirements.

\section{Conclusion}

We presented TrajScale, a decoupled mixed-resolution sampling framework
that combines a low-resolution prefix with a short high-resolution suffix.
EAA calibrates the handoff prediction toward the converged low-resolution
endpoint, CRLU upsamples clean latents across spatial scales, and TRR
reconstructs a schedule-consistent high-resolution state. All trainable
components use supervision generated by the frozen base model.

TrajScale-Q and TrajScale-E achieve \(1.83\times\) and \(2.22\times\) speedups
while retaining \(97.76\%\) and \(97.53\%\) of Native-HR VBench-5
performance. CRLU consistently outperforms trilinear interpolation,
while EAA improves endpoint alignment and localized visual detail across
standard and distilled samplers.
